\documentclass{momentum}
\usepackage{unicode-math}
\setmathfont{Garamond-Math.otf}
\XeTeXgenerateactualtext=1
\setstretch{1.13}

\title{Critical points and sphere insertions}
\author{Raeez Lorgat}
\date{May 27, 2025}
\shorttitle{Critical points and sphere insertions}
\hypersetup{
  pdftitle={Critical points and sphere insertions},
  pdfauthor={Raeez Lorgat},
  pdfsubject={Closed amplitudes and marked reconstruction in Landau–Ginzburg theory},
  pdflang={en}
}

\usepackage{etoolbox}
\patchcmd{\maketitle}
  {\thetitle\par}
  {\thetitle\par\vspace{0.25em}%
   {\fontsize{15pt}{18pt}\selectfont Closed amplitudes in a Landau--Ginzburg family\par}}
  {}{\PackageError{sphere-print}{Missing title insertion point}{}}
\pretocmd{\thebibliography}{\phantomsection\addcontentsline{toc}{section}{\refname}}
  {}{\PackageError{sphere-print}{Missing bibliography insertion point}{}}
\apptocmd{\thebibliography}{\raggedright}
  {}{\PackageError{sphere-print}{Missing bibliography formatting point}{}}
\newcommand{\Tr}{\operatorname{Tr}}
\newcommand{\Res}{\operatorname{Res}}

\begin{document}
\maketitle

\noindent
Take one complex field \(x\) with potential \(W_t(x)=x^3/3-tx\), where
\(t\) is a complex parameter. A critical point is a solution of
\(W_t'(x)=0\). Here the equation is \(x^2=t\). For \(t\ne0\) the two
solutions are \(s\) and \(-s\), with \(s^2=t\). They meet at the origin
when \(t=0\). Can a correlation function on the sphere follow the two
solutions through their collision?

In the B-twisted Landau--Ginzburg theory a supersymmetry acts on local
fields by a differential \(Q_{\mathrm B}\) with \(Q_{\mathrm B}^2=0\).
Fields annihilated by \(Q_{\mathrm B}\), taken modulo fields of the form \(Q_{\mathrm B}a\),
give the protected local observables. A polynomial \(f(x)\) represents
such an observable and can be inserted at a point of the sphere.
Fix the holomorphic volume form \(dx\). The sphere correlation function
for this insertion is the sum of the residues of \(f(x)\,dx/W_t'(x)\)
at the critical points \cite{vafa}.
A residue is the coefficient of \(du/u\) in a local coordinate \(u\)
vanishing at the point. At a simple zero \(\alpha\) of \(W_t'\) it is
\(f(\alpha)/W_t''(\alpha)\). For the cubic the two contributions are
\(f(s)/(2s)\) and \(-f(-s)/(2s)\). Inserting \(1\) gives zero, inserting
\(x\) gives one, and inserting \(x^3\) gives \(s^2=t\).
As \(s\) tends to zero, the two vacuum weights \(1/(2s)\) and \(-1/(2s)\)
diverge separately. With \(f\) inserted, their contributions add to
\((f(s)-f(-s))/(2s)\), a divided difference with limit \(f'(0)\).
At the collision this is precisely the residue of \(f(x)\,dx/x^2\).
In particular \(x^3\) records the splitting parameter right up to the
collision. What if the midpoint of the roots can also move?

\section{Reconstructing the marked potential}

Allow the linear as well as the constant coefficient of the critical
equation to vary. Write it as
\(P(x)=x^2+ax+b=0\). By keeping \(x\) marked we mean that the same
polynomial in the same coordinate will be inserted for each choice of
\(a,b\). Write \(\tau(f)\) for its sphere value, still with volume \(dx\).
At distinct roots \(\alpha,\beta\) the two residues give
\[
  \tau(f)=\frac{f(\alpha)-f(\beta)}{\alpha-\beta}.
\]
In particular \(\tau(1)=0\) and \(\tau(x)=1\). Adding a multiple of \(P\)
to \(f\) adds a polynomial differential to \(f\,dx/P\). The added
differential is regular at both roots and has zero residue there. We can therefore use \(x^2=-ax-b\) inside any
sphere insertion. Set \(M_r=\tau(x^r)\). Reduction first of \(x^2\) and
then of \(x^3=-ax^2-bx\) gives
\[
 M_2=-a,\qquad M_3=a^2-b,\qquad
 a=-M_2,\qquad b=M_2^2-M_3 .
\]
For roots \(1,3\) the weights are \(-1/2,1/2\). The two readings are
\(4,13\), so the recovered polynomial is \(x^2-4x+3\). At a double root \(r\), the form is \(f(x)\,dx/(x-r)^2\), with residue
\(f'(r)\). Thus \(M_2=2r\) and \(M_3=3r^2\). The same two readings
recover \(a=-2r\) and \(b=r^2\), the coefficients of \((x-r)^2\).

The reduction of powers used only that the leading coefficient of the
critical polynomial was one. Division by such a monic polynomial makes
sense with coefficients in any commutative \(\mathbb Q\)-algebra \(R\),
even when a root cannot be chosen in \(R\). Let
\(P(x)=x^d+a_{d-1}x^{d-1}+\cdots+a_0\), \(d\ge1\), be the derivative
of a potential \(W\), with \(W(0)=0\). The closed fields form the
\emph{Jacobi algebra} \(A=R[x]/(P)\). Its elements have unique
representatives of degree less than \(d\). Define \(\tau(f)\) as the
coefficient of \(x^{d-1}\) in this remainder. To see the residue in this
definition, expand \(f/P\) in descending powers of \(x\). Its polynomial
part has no \(x^{-1}\) term, and dividing the remainder by the monic
denominator makes its \(x^{d-1}\) coefficient the coefficient of
\(x^{-1}\) in \(f/P\). Over \(\mathbb C\) this equals the sum of the
finite residues of \(f\,dx/P\), including those at multiple roots.

The moments \(M_0,\ldots,M_{d-2}\) vanish and \(M_{d-1}=1\).
For \(j\ge0\) the relation \(x^jP(x)=0\) in \(A\) gives
\(M_{d+j}+a_{d-1}M_{d+j-1}+\cdots+a_0M_j=0\).
At \(j=0\) only two terms survive, so \(a_{d-1}=-M_d\).
For \(d\ge2\) the next equation gives
\(a_{d-2}=-M_{d+1}-a_{d-1}M_d\).
At each step the new coefficient multiplies \(M_{d-1}=1\), and all
coefficients of higher powers of \(x\) are already known. Thus the
\(d\) fixed insertions \(x^d,\ldots,x^{2d-1}\) recover \(P\) over \(R\).
The same recurrence determines every later moment. Collecting them in
a formal variable \(z\) gives
\begin{equation}
 \sum_{r\ge0} M_r z^r
  =\frac{z^{d-1}}{1+a_{d-1}z+\cdots+a_0z^d}.
 \label{eq:moments}
\end{equation}
Integrating the recovered \(P\) gives \(W\) with the chosen value
\(W(0)=0\). Adding a constant to \(W\) leaves all the sphere readings
unchanged, since the residue uses its derivative.

Could a smaller set of insertions detect the same coefficients? Take
a characteristic-zero field \(k\) and vary \(x^d\) to
\(x^d+\epsilon q(x)\), where \(\deg q<d\) and \(\epsilon^2=0\).
The coefficient of \(\epsilon\) records a first-order variation.
Inverting the denominator gives \(x^{-d}-\epsilon q(x)x^{-2d}\).
Thus the change in the sphere reading is
\begin{equation}
 \tau_{x^d+\epsilon q}(f)-\tau_{x^d}(f)
   =-\epsilon [x^{2d-1}]\bigl(f(x)q(x)\bigr),
 \label{eq:tangent}
\end{equation}
where brackets extract the indicated coefficient. Each fixed polynomial
insertion gives one linear functional on the \(d\) coefficients of \(q\).
Fewer than \(d\) such functionals have a common nonzero kernel. For fixed
monomial insertions one can say exactly which readings are indispensable.
If the power \(x^j\) is missing, with \(d\le j\le2d-1\), take
\(q=x^{2d-1-j}\). Among all powers of \(x\), only \(x^j\) detects this
variation. Thus every one of \(x^d,\ldots,x^{2d-1}\) is necessary for a
reconstruction by fixed monomial insertions over all parameter rings.

\section{The diagonal and the handle}

The readings now determine the entire sphere pairing. Two insertions
\(f,h\) have the joint value \(\tau(fh)\). How does this pairing tell us
what to insert for a handle? Cutting the handle open produces two
boundary circles. To sew them back together one sums a basis of states
on one circle against the dual basis on the other. At our first double
point, \(P=x^2\), the sphere gives \(\tau(1)=\tau(x^2)=0\) and
\(\tau(x)=1\). Its pairing on \(1,x\) has matrix
\(\left(\begin{smallmatrix}0&1\\1&0\end{smallmatrix}\right)\).
The dual to \(1\) is \(x\), and the dual to \(x\) is \(1\). Sewing
therefore places \(1\otimes x+x\otimes1\) on the two circles.
Writing \(x,y\) for the coordinates on the first and second factors,
this is \(x+y=(x^2-y^2)/(x-y)\).

For general \(P\) the same divided difference defines the diagonal class
\[
 \mathcal D(x,y)=\frac{P(x)-P(y)}{x-y}\quad\text{in } A\otimes_R A .
\]
First divide in the polynomial ring \(R[x,y]\), where the numerator is
a multiple of \(x-y\). Then impose \(P(x)=P(y)=0\). In \(A\otimes_R A\)
it obeys \((x-y)\mathcal D=0\), so \(f(x)\mathcal D=f(y)\mathcal D\).
Its coefficient of \(x^{d-1}\) is one. Applying the sphere functional
in the first variable consequently gives
\((\tau_x\otimes1)(f(x)\mathcal D)=f(y)\), the defining property of
the inverse pairing. In the power basis the same conclusion follows
from \(\tau(x^{i+j})=0\) for \(i+j<d-1\) and
\(\tau(x^{i+j})=1\) for \(i+j=d-1\). Reversing the column order gives
a triangular matrix with unit diagonal.

Joining the two circles multiplies their labels in \(A\), which sets
\(y=x\) and sends \(\mathcal D\) to \(e=P'(x)\). Thus a handle inserts
one copy of \(e\). Write \(A_g(f)\) for the closed amplitude on a
connected surface with \(g\) handles, where \(f\) is the product of
its marked fields. Write \(m_h\) for multiplication by \(h\) in \(A\).
Repeated sewing gives
\begin{equation}
 A_g(f)=\tau(fe^g),\qquad
 A_{g+1}(f)=\Tr_{A/R}(m_fm_e^g).
 \label{eq:handle}
\end{equation}
Indeed, if \(u_i,v_i\) are dual bases, then
\(e=\sum_i u_i v_i\) and \(\Tr(m_h)=\sum_i\tau(v_i h u_i)=\tau(he)\).
In particular the torus with no insertions has amplitude
\(\Tr_{A/R}(m_1)=d\).
What if the potential is multiplied by a scalar while \(dx\) stays fixed?
For \(W=x^D\), \(D\ge2\), the denominator in the residue is
\(D x^{D-1}\). Relative to \(W=x^D/D\) this divides \(\tau\) by \(D\)
and multiplies its inverse pairing and \(e\) by \(D\). Their product
in the torus amplitude still gives \(D-1\).

The higher genera can now be obtained from a single finite matrix.
As a formal series, \(\sum_{g\ge0}A_g(f)z^g=\tau((1-zm_e)^{-1}f)\).
The adjugate formula gives a denominator \(\det(1-zm_e)\) of degree
at most \(d\) with constant term one. Reduction by monic \(P\), extraction
of one coefficient and multiplication in \(A\) use only ring operations.
The sewing formulas therefore commute with every change \(R\to R'\)
of parameter ring. This includes specializing two distinct roots to a
double root.

\section{The fibre at a collision}

Consider a fibre over an algebraically closed field \(k\) of characteristic zero.
The local part of \(A\) at a root \(\alpha\) of multiplicity \(m\) is
\(k[u]/(u^m)\), with \(u=x-\alpha\). Factoring off this multiplicity gives
\(P(x)=(x-\alpha)^m V(x)\) with \(V(\alpha)\ne0\). On the sphere the contribution of this root
is the coefficient of \(u^{m-1}\) in \(f(\alpha+u)/V(\alpha+u)\).
On a surface of positive genus formula \eqref{eq:handle} gives the
ordinary trace of multiplication by \(h=f(P')^{g-1}\) on this local
algebra. In the basis \(1,u,\ldots,u^{m-1}\) its matrix is triangular
with the constant \(h(\alpha)\) on each diagonal entry. A simple root
therefore contributes \(f(\alpha)P'(\alpha)^{g-1}\), also at \(g=0\)
by the residue formula. A multiple root contributes \(m f(\alpha)\)
at \(g=1\). For \(g\ge2\) its contribution vanishes because \(P'(\alpha)=0\).

How does that vanishing look before specialization? Over \(k[t]\) take
\(W_t=x^{d+1}/(d+1)-tx\) with \(d\ge2\). Now \(x^d=t\) and \(e=d x^{d-1}\). For \(r,g\ge0\),
the sphere insertion corresponding to \(x^r\) on the surface with
\(g\) handles is \(x^r e^g=d^g x^{r+(d-1)g}\). Each group of \(d\)
powers of \(x\) contributes one factor of \(t\). Exactly \(d-1\)
powers of \(x\) have to remain for the sphere trace to survive.
The number of groups must therefore be
\(h=(r+(d-1)g-(d-1))/d\). This gives
\begin{equation}
 A_g(x^r)=
 \begin{cases}
 d^g t^h,&h\in\mathbb Z_{\ge0},\\
 0,&\text{otherwise}.
 \end{cases}
 \label{eq:family}
\end{equation}
For the vacuum \(r=0\), the condition on \(h\) singles out \(g=1+jd\)
with \(j\ge0\). Summing these terms gives
\(\sum_{g\ge0} A_g(1)z^g=dz/(1-d^d t^{d-1}z^d)\).
For the cubic that opened the calculation, \(A_3(1)=8t\).
Set \(t=\epsilon\) with \(\epsilon^2=0\). The underlying point
has \(t=0\), but the amplitude over the parameter ring is \(8\epsilon\).

For a sum of potentials in independent coordinates \(x_i\), use the
ordered product volume form. The quotient by their separate critical
equations is the tensor product of the one-variable algebras. The residue
trace and its inverse pairing factor, so the amplitude for
\(\prod_i x_i^{r_i}\) is the product of \eqref{eq:family} for
the individual \(d_i,t_i\), including when only some \(t_i\) vanish.

Which critical points can support a boundary condition at a specified
value of the potential? The sphere
used \(W'\), whereas a boundary at value \(w\) depends on \(W-w\).
Return to the critical roots \(1,3\). With \(W(0)=0\) their potential is
\(W=x(x-3)^2/3\), so \(W(1)=4/3\) and \(W(3)=0\). Both points contribute
to the closed torus amplitude of two. A boundary at value zero can
see a critical point on its own fibre only at \(3\).

In this model a finite free boundary condition at value \(w\) is a
matrix factorization. It consists of a finite free module with an
even and an odd part, and an odd matrix \(Q\) interchanging the two
parts such that \(Q^2=(W-w)1\). At \(x=3\), write \(u=x-3\) and
complete at \(u=0\), using the formal power series ring \(k[[u]]\).
Take one copy of that ring in each parity. The two maps in \(Q\) are
multiplication by \(u\) and by \(u(u+3)/3\). Their product is \(W\).
The identity of the boundary would be homotopic to zero if some odd
matrix \(H\) satisfied \(QH+HQ=1\). Both maps in \(Q\) vanish at \(u=0\),
so this equality would give \(0=1\) there. The boundary survives.
Over \(W=0\) the first map presents the residue field \(k[[u]]/(u)\).
Alternating the two maps continues its resolution. This factorization,
called the stabilization of the residue field, generates the local
category \cite[Thm.~4.1]{dyck}.

What happens to such a boundary when its fibre contains no critical
point? For the family in \eqref{eq:family} over the rational function field
\(k(t)\), every critical root satisfies \(\alpha^d=t\). Its critical
value is \(-dt\alpha/(d+1)\), which is nonzero. Thus \(W_t\) and \(W_t'\)
have no common root. At value zero choose polynomials \(b,c\) with
\(bW_t+cW_t'=1\). Differentiating \(Q^2=W_t1\) gives
\(Q\,\partial_xQ+(\partial_xQ)Q=W_t'1\). Hence
\[
 H=\tfrac12 bQ+c\,\partial_xQ,\qquad QH+HQ=1 .
\]
Every such boundary factorization is contractible on this generic
fibre, while the closed torus amplitude is \(d\).

Return to the boundary that survives at \(3\). Its local ring of
closed fields is \(k[[u]]/(u(u+2))=k\), since \(u+2\) is a unit.
One sees this ring on the boundary side by putting the same theory on
both sides of a line and letting it act by the identity. Use the completed local coefficient ring \(k[[x-3,y-3]]\), so that both
field coordinates lie near \(3\). Put \(v=(W(x)-W(y))/(x-y)\), a
polynomial viewed in this ring. The two maps
\(x-y\) and \(v\) multiply to \(W(x)-W(y)\) and form the matrix
factorization for this identity defect. Its even endomorphisms in
cohomology are scalars modulo \(x-y\) and \(v\). An even diagonal
endomorphism must have equal entries to commute with the two maps,
and an odd homotopy changes that common entry by a combination
of the two maps. Setting \(y=x\) gives \(v(x,x)=W'(x)\) and hence
the local Jacobi algebra \cite[Prop.~6.3 and Cor.~6.4--6.5]{dyck}.
The calculation identifies the cohomology ring of the defect. On a disk,
one also needs a trace on boundary morphisms, given by the Kapustin--Li
residue \cite[Thm.~3.4 and \S5.1, Eq.~(33)]{dm}. The local calculation
accounts for the vacuum at \(3\). To recover the closed torus amplitude
of two, one must also include the sector at critical value \(4/3\),
supported at \(1\).

\begin{thebibliography}{9}\small
\bibitem{vafa} C.~Vafa, \emph{Topological Landau--Ginzburg models},
Modern Phys.\ Lett.\ A \textbf{6} (1991), 337--346.
\href{https://doi.org/10.1142/S0217732391000324}{doi:10.1142/S0217732391000324}.
\bibitem{dyck} T.~Dyckerhoff, \emph{Compact generators in categories of
matrix factorizations}, Duke Math.\ J.\ \textbf{159} (2011), 223--274.
\href{https://arxiv.org/abs/0904.4713}{arXiv:0904.4713}.
\bibitem{dm} T.~Dyckerhoff and D.~Murfet, \emph{The
Kapustin--Li formula revisited}, Adv.\ Math.\ \textbf{231} (2012),
1858--1885. \href{https://arxiv.org/abs/1004.0687}{arXiv:1004.0687}.
\end{thebibliography}
\end{document}
